\section*{Chapter \ref{ch:qm:point-processes-and-hawkes-processes} --- Point Processes and Hawkes Processes}

\begin{solution}{exo:qm:point-processes-and-hawkes-processes:1}
$n = 0.7$; $\bar\lambda = 0.1/0.3 = 0.333$ per second; $0.333 \times 23\,400 = 7\,800$ events.
\end{solution}

\begin{solution}{exo:qm:point-processes-and-hawkes-processes:2}
$1/(1 - 0.7) = 3.33$ events per immigrant; immigrants are $1 - n = 30\%$ of events.
\end{solution}

\begin{solution}{exo:qm:point-processes-and-hawkes-processes:3}
$\Lambda(t) = at + bt^2/2$; $\P(N_T = 0) = e^{-aT - bT^2/2}$.
\end{solution}

\begin{solution}{exo:qm:point-processes-and-hawkes-processes:4}
$\ln 2/\beta = 0.69$ seconds; the jump $\alpha = 0.7$ is seven times the baseline $0.1$.
\end{solution}

\begin{solution}{exo:qm:point-processes-and-hawkes-processes:5}
With $\gamma = \beta(1 - n) = 0.3$ and $n(2 - n)\beta/(2(1 - n)) = 1.517$: $1 + 2 \times 1.517\,(1/0.3 - (1 - e^{-0.3})/0.09) = 2.38$ at
one second, 10.5 at a minute, and $1/0.3^2 = 11.1$ in the limit.
\end{solution}

\begin{solution}{exo:qm:point-processes-and-hawkes-processes:6}
Mean 2, variance $\E[\Lambda] + \Var(\Lambda) = 2 + 1 = 3$. It is a mixed Poisson (Cox) process: overdispersed without any
self-excitation, which is why overdispersion alone does not identify a \omterm{def:qm:point-processes-and-hawkes-processes:hawkes}{Hawkes process}.
\end{solution}

\begin{solution}{exo:qm:point-processes-and-hawkes-processes:7}
Mean 0.700, standard deviation 0.011 over forty days of about 7\,800 events.
\end{solution}

\begin{solution}{exo:qm:point-processes-and-hawkes-processes:8}
A constant baseline cannot follow the intraday pattern of activity, and the fit explains busy periods by
excitation with a long kernel: a day with no excitation at all but a U-shaped rate gives 0.87. Model the
baseline's variation (or rescale time by the intraday profile) before reading the \omterm{def:qm:point-processes-and-hawkes-processes:hawkes}{branching ratio}, and check
the residuals by time of day.
\end{solution}

\begin{solution}{pb:qm:point-processes-and-hawkes-processes:1}
\textbf{1.} $\alpha/\beta = 0.7$.
\textbf{2.} $0.333$ per second; 7\,800 trades.
\textbf{3.} $3.33$; 30\% immigrants.
\textbf{4.} $\ln 2 = 0.69$ seconds.
\textbf{5.} 2.38 and 10.5.
\textbf{6.} 8\,006.
\textbf{7.} $\hat\mu = 0.103$, $\hat\alpha = 0.696$, $\hat\beta = 0.995$: $\hat n = 0.700$.
\textbf{8.} $0.700 \pm 0.011$.
\textbf{9.} Hawkes fit: mean 1.000, variance 0.98. Poisson fit: mean one by construction, variance 4.63.
\textbf{10.} 69.9\% (7\,498 events).
\textbf{11.} Three times higher at the open and the close; mean rate $0.333$ per second, $0.2$ at midday.
\textbf{12.} 7\,730.
\textbf{13.} $\hat n = 0.87$ with $\hat\beta = 0.0076$ per second: a kernel lasting about two minutes.
\textbf{14.} Mean 0.86 over twenty seasonal days; 0.04 over twenty constant-rate Poisson days.
\textbf{15.} A slowly varying rate makes events cluster in time; with a constant baseline, the only way the model
can put more events where there are more is through excitation with a long kernel.
\textbf{16.} Estimate a time-varying baseline $\mu(t)$ jointly with the kernel, or rescale time by the estimated
intraday profile first; check residuals separately at the open, midday and close.
\textbf{17.} That the \omterm{def:qm:point-processes-and-hawkes-processes:hawkes}{branching ratio} depends on the kernel's form, the window, the baseline and outliers:
it is a fragile number, and comparisons across periods need the same method.
\textbf{18.} To forecast short-term activity and adverse flow (widen or pull quotes when the intensity spikes),
to schedule execution away from bursts, and to generate realistic order flow for backtests.
\textbf{19.} Named result: \emph{the self-exciting tape}: on days with a true \omterm{def:qm:point-processes-and-hawkes-processes:hawkes}{branching ratio} of 0.7 the fit
recovers $0.700 \pm 0.011$; on days with no excitation but a U-shaped baseline it reports 0.87 (0.86 on average),
against 0.04 for constant-rate Poisson days.
\textbf{20.} Only when the baseline's own variations are modelled and the time-rescaled residuals pass.
\end{solution}

\begin{solution}{iq:qm:point-processes-and-hawkes-processes:1}
A \omterm{def:qm:point-processes-and-hawkes-processes:pp}{counting process} whose intensity is a baseline plus a decaying kernel summed over past events, so each event
raises the rate of further events. The \omterm{def:qm:point-processes-and-hawkes-processes:hawkes}{branching ratio}, the integral of the kernel, is the mean number of events
each event triggers, and the long-run share of activity that is endogenous.
\iqlookfor{the intensity formula and the two readings of the \omterm{def:qm:point-processes-and-hawkes-processes:hawkes}{branching ratio}.}
\end{solution}

\begin{solution}{iq:qm:point-processes-and-hawkes-processes:2}
Take expectations in stationarity: $\bar\lambda = \mu + n\bar\lambda$, so $\bar\lambda = \mu/(1 - n)$. At $n \ge 1$ each event produces on
average at least one more, clusters are infinite with positive probability, and no stationary version
exists: the process explodes.
\iqlookfor{the fixed-point argument and criticality.}
\end{solution}

\begin{solution}{iq:qm:point-processes-and-hawkes-processes:3}
By \omterm{def:qm:point-processes-and-hawkes-processes:thinning}{thinning}: the exponential intensity decays between events, so the current intensity bounds it until the
next event; propose at that rate and accept with probability intensity over bound. Or generate immigrants,
then Poisson numbers of children generation by generation.
\iqlookfor{Ogata's \omterm{def:qm:point-processes-and-hawkes-processes:thinning}{thinning} or the cluster representation, and why both are exact.}
\end{solution}

\begin{solution}{iq:qm:point-processes-and-hawkes-processes:4}
The intensity at $t_i$ needs $A_i = \sum_{j<i}e^{-\beta(t_i - t_j)}$, and $A_i = e^{-\beta(t_i - t_{i-1})}(1 + A_{i-1})$; the integral
of the intensity is a sum of one term per event. Both are one pass over the data.
\iqlookfor{the recursion.}
\end{solution}

\begin{solution}{iq:qm:point-processes-and-hawkes-processes:5}
As a short-horizon forecast of how busy the next seconds will be and on which side: widen or skew quotes when the
opposite side's intensity spikes, since bursts carry adverse selection; size inventory limits by expected
flow; and in the simulator, generate clustered flow instead of Poisson flow.
\iqlookfor{intensity as a forecast of flow and of adverse selection.}
\end{solution}

\begin{solution}{iq:qm:point-processes-and-hawkes-processes:6}
Rescale time by the fitted \omterm{def:qm:point-processes-and-hawkes-processes:intensity}{compensator} and test that the gaps are independent unit exponentials: mean and
variance, a QQ plot, a Kolmogorov--Smirnov test (chapter~12), autocorrelation of the gaps, and all of it by time
of day and out of sample.
\iqlookfor{the time-rescaling theorem as a goodness-of-fit tool.}
\end{solution}
